%%%%%
%
%% Dans ce fichier, les �ventuels accents sont cod�s
%% en ISO-8859-15 (aka latin9)
%% La ligne suivante donne trois exemples 
%% e accent aigu : �, e accent grave : �, e-dans-l'o : 1/2

\section[Humidit� relative � partir du point de ros�e ou de givre,
seuill�e � 102\%]{Humidit� relative � partir de la temp�rature
du point de ros�e ou de givre, seuill�e � 102\%}

\itemstyle{Op�rateur~:} \operfmt{HumW\_T\_Tdf}
\label{OP-HumW_T_Tdf}
\index{HumW\_T\_Tdf}

%%%%%OBJET%%%%%
\itemstyle{ Objet~: }
Cet op�rateur calcule l'humidit� relative � l'eau liquide � partir de la
temp�rature 
du point de ros�e ou de givre et de la temp�rature de l'air.

%%%%%ARGUMENTS%%%%
\bigskip
%% Dans ce fichier, les �ventuels accents sont cod�s
%% en ISO-8859-15 (aka latin9)
%% La ligne suivante donne trois exemples 
%% e accent aigu : �, e accent grave : �, e-dans-l'o : 1/2
%%%%%ENTREES%%%%
\begin{tabarg}
\multicolumn{3}{l}{\titretab {Entr�es}}\\
\hline
$T_s$ & Vecteur, \texttt{float} & Temp�rature de l'air ($K$ ou $^\circ C$) \\
$T_{df}$ & Vecteur, \texttt{float} & Temp�rature du point de ros�e ou de givre ($K$ ou $^\circ C$) 
\end{tabarg}

\medskip
%%%%%SORTIES%%%%
\begin{tabarg}
\multicolumn{3}{l}{\titretab {Sorties}}\\
\hline
$H_u$ & Vecteur, \texttt{float} & Humidit� relative � l'eau liquide (\%)\\
\end{tabarg}


%%%%%FORMULE%%%%
\itemstyle{Formule}

%% Formule de l'operateur : HumW_T_Tdf
%% Dans ce fichier, les �ventuels accents sont cod�s
%% en ISO-8859-15 (aka latin9)
%% La ligne suivante donne trois exemples 
%% e accent aigu : �, e accent grave : �, e-dans-l'o : �
\[
e = 
\begin{cases}
 e_w(T_{df}) \text{ si } T_{df} \ge 0.0 \text{�C}\\
 e_{wi}(T_{df}) \text{ sinon}\\
\end{cases}
H_u = 100 \times \frac{e}{e_w(T_s)} 
\]
Si $H_u > 102 \%$, $Hu = 103 \%$.


%%%%COMMENTAIRE%%%%
\itemstyle{Commentaire~:}

Les fonctions $ew()$ et  $ewi()$ sont d\'ecrites en annexe
\ref{annexe-thermo}.

%%%%%AUTEUR%%%%
\itemstyle{Auteur~:}
CNRM/GMEI/TRAMM

%%%%%REFERENCE%%%%
\itemstyle{R�f�rences~:}

\emph{M�t�orologie G�n�rale}, Triplet-Roche, �dition de 1971\cite{Triplet-Roche}, p59.
